\documentclass{article}
\usepackage[utf8]{inputenc}
\usepackage{booktabs}
\usepackage{hyperref}
\usepackage{geometry}

\geometry{a4paper, margin=2.5cm}

\title{\textbf{Radiografia e Manual de Integração da API do Moodle (FICV)}}
\author{Plataforma ARKOS}
\date{\today}

\begin{document}

\maketitle

\section{Introdução e Visão Geral}
Este documento apresenta o mapeamento definitivo das variáveis estruturadas, os desafios de dados não estruturados, e os processos de ETL (Extract, Transform, Load) necessários para garantir a máxima usabilidade da integração Moodle com o ecossistema ARKOS, visando a construção de séries temporais e suporte à decisão gerencial acadêmica.

\section{Dados Estruturados e Extração Direta}
Os seguintes campos são obtidos de fontes relacionais nativas do Moodle e estão PRONTOS para persistência direta em tabelas como \texttt{dados\_moodle\_cursos} no Supabase:

\vspace{0.3cm}
\begin{table}[htbp]
  \centering
  \caption{Mapeamento de Variáveis Estruturadas}
    \begin{tabular}{lll}
    \toprule
    \textbf{Variável} & \textbf{Tipo de Dado} & \textbf{Origem Moodle (API / DB)} \\
    \midrule
    \texttt{id} & Integer & Core User Index \\
    \texttt{fullname} & String & \texttt{core\_enrol\_get\_enrolled\_users} \\
    \texttt{email} & String (Email) & \texttt{core\_enrol\_get\_enrolled\_users} \\
    \texttt{cpf} & String (Regex) & \texttt{customfields} (EAD Profiles) \\
    \texttt{semestre} & String (Ex: 2024.1) & \texttt{customfields} (EAD Profiles) \\
    \texttt{media\_geral} & Float & Gradebook Aggregate (\texttt{graderaw}) \\
    \bottomrule
    \end{tabular}
\end{table}

\section{Dados Não Estruturados e Processo de ETL}
Algumas variáveis cruciais vêm em formatos híbridos (strings mistas ou dados aninhados), exigindo tratamentos específicos de \textbf{Regex} e \textbf{Splitting}:

\begin{itemize}
    \item \textbf{Notas Descritivas (\texttt{notas\_disciplinas})}:
    \begin{itemize}
        \item \textit{Formato Atual}: String composta (Ex: ``Módulo 1: 9.5 | Módulo 2: 7.0'').
        \item \textit{Ajuste ETL}: Fazer o split pelo caractere ``|'' e utilizar Expressor Regular para extrair [Nome\_Modulo, Nota] em um \textit{JSONB} ou tabela auxiliar de histórico mensal.
    \end{itemize}
    \item \textbf{Papéis (Professor vs Aluno)}:
    \begin{itemize}
        \item \textit{Desafio}: A API de incritos traz todos os cargos em uma lista flat.
        \item \textit{Ajuste ETL}: É fundamental filtrar os registros na camada de serviço (Back-End) onde \texttt{roles.shortname == 'editingteacher'} para preencher a variável de DOCENTE.
    \end{itemize}
\end{itemize}

\section{Geração de Séries Temporais para Decisão}
Com o ETL estruturado, os dados podem ser utilizados para previsões analíticas:

\begin{enumerate}
    \item \textbf{Análise de Churn / Evasão}: Alunos que constam como ``Em Curso'' mas não possuem inserção de novas notas nos últimos 45 dias geram Score de Risco para a instituição.
    \item \textbf{Janela Histórica}: O Moodle não impõe limitação de datas na listagem. Dados cobrem todo o histórico FICV, permitindo agrupar aprovações por Semestre e Categoria.
    \item \textbf{Docência Criativa}: Correlacionar a média geral de turmas à variáveis como Titulação do professor (cadastrada em custom fields) para avaliar impacto do docente na taxa de aprovação.
\end{enumerate}

\section{Conclusão}
A integração está estruturada para consumo analítico. Recomenda-se a ativação de um Cron automatizado para povoar o Supabase de forma incremental a cada hora, sustentando painéis de BI em tempo real.

\end{document}
